\section{Introduction}
A gas-concentration estimate depends both on the sensor response and on the reference information used to interpret that response. Recalibration can therefore change performance without changing the sensor hardware or the upstream prediction model. For a reference-assisted method, a favorable comparison can have at least two explanations: the method uses existing calibration data more effectively, or it gains useful information from newly requested reference measurements. Distinguishing these explanations is essential when a method is presented as reducing calibration effort.

This distinction is particularly relevant to metal-oxide (MOX) sensor arrays. Device-to-device variation and environmental sensitivity make calibration a substantive part of the measurement procedure, rather than only a preprocessing step. Work on global calibration has examined how models built from several sensor replicas transfer to another unit \cite{miquel2022}. Calibration experiments using randomized gas mixtures have also demonstrated that the conditions represented in the calibration set affect subsequent prediction performance \cite{baur2020}. These studies motivate explicit accounting of what was calibrated, which reference information was available, and which experimental population a result represents.

Additional reference access raises a different question. A sensor system may already possess a substantial historical calibration set, while obtaining new reference values is limited. A reference-checking policy must then be judged against an estimator that makes effective use of that same historical set. A comparison only with an uncorrected or weakly corrected baseline can credit a reference policy for an advantage already attainable before any new check. Conversely, a failure of one update rule to benefit from additional labels does not prove that reference measurements are generally uninformative.

This paper addresses these issues through a controlled comparative evaluation of CO concentration estimates. It uses the public UCI temperature-modulated gas-sensor dataset \cite{uci487}, with complete chamber exposures kept together and model fitting, historical correction calibration, and assessment separated chronologically. Five point-estimation rules consume the same reference-query masks. This separates estimator effects from schedule effects while charging one reference per shared chamber exposure. The primary endpoint is the fraction of pre-query reference-relative errors exceeding a fixed 1 ppm research tolerance; mean absolute error (MAE), other tolerances, and interval diagnostics provide complementary views.

The evaluation has two distinct stages: chronological development on D6/D8/D10 and a completed, separately frozen fixed-D10 temporal confirmation on D11/D12/D13 (Section~\ref{sec:c5}). The frozen descriptive benefit criterion requires lower strict 1 ppm exceedance and nonpositive MAE change pooled and separately in both families. Under that criterion, historical benefit is heterogeneous, whereas the prespecified incremental affine envelope 20-check comparison meets the confirmation aggregate criterion. Pooled confirmation MAE changes from 1.3553 to 1.3257 ppm and strict exceedance from 43.06\% to 42.14\%, reversing the adverse development B20 MAE pattern. Adverse session/family cells and unsupported envelope-schedule superiority remain part of the finding.

Three contributions follow: an information-matched chronological protocol spanning development and temporal confirmation, a reporting identity separating historical correction from incremental reference access, and a diagnostic distinction between containment and useful assurance. The comparison uses established correction algorithms. Forgetting and one native schedule are declared development extensions retained in the frozen confirmation grid. Both stages use the same apparatus; temporal confirmation does not establish independent-array generalization.

\section{Related work and scope of the comparison}\label{sec:related}
\paragraph{Calibration information and sensor transfer.}
Miquel-Ibarz et al.~\cite{miquel2022} evaluated individual and global calibration models for temperature-modulated MOX sensors. Their study uses FIS replicas and examines calibration cost, transfer to unseen units, and temporal behavior. Its processed data description, concentration subset, and campaign structure differ from the 14-channel, 13-file release used here. We therefore do not identify the two evaluations as identical-data replications or compare their numerical scores as interchangeable. The present question concerns a calibrated array and reference access during assessment, not removal of the need for an initial calibration of a new unit.

Baur et al.~\cite{baur2020} studied calibration with randomized mixtures rather than a conventional sequence of isolated gas concentrations. Their work emphasizes the representativeness of the initial calibration experiment. Here the archived calibration conditions are fixed. The question is how much additional value a limited set of subsequently revealed reference labels contributes after those historical conditions have already been used.

\paragraph{Recalibration and reference-assisted measurement.}
G\"abel and Hertig~\cite{gaebel2026} evaluated repeated recalibration of low-cost ozone and particulate-matter sensors using a year-long reference collocation campaign. Their findings support the importance of calibration age and environmental conditions in their field setting. Our short chamber replay does not replace such a field experiment: its label-access budgets, cumulative updates, and day resets define a different measurement setting. Gruber et al.~\cite{gruber2023} considered reference-assisted co-calibration in distributed homogeneous sensor networks, including interpolation and propagation of reference uncertainty. In contrast, the present released CO values do not supply an exposure-level hard uncertainty bound or the joint reference-error structure needed for a traceable uncertainty budget.

\paragraph{Bounded-error estimation.}
Set-membership identification and parameter bounds from bounded-error observations have a longstanding literature \cite{walter1990}. The max/min envelope, midpoint prediction, empirical slope, and two-endpoint noise allowance used in this study have direct antecedents in Lipschitz interpolation \cite{calliess2017}. Sequential Lipschitz optimization also distinguishes assumed regularity from the task of estimating it \cite{malherbe2017}. These precedents do not validate the empirical temporal residual bound used here. The established constructions are retained as comparators and inspected for validity, rather than presented as newly derived measurement guarantees.

The comparison therefore concerns the interaction of initial reference use, later reference access, point correction, and reporting criteria. It does not reproduce a published active-learning algorithm, establish an exhaustive ranking of recalibration methods, or demonstrate a complete uncertainty-aware co-calibration chain.

\paragraph{Drift compensation, field calibration and recursive identification.}
Classifier ensembles have been studied for long-term MOX gas discrimination \cite{vergara2012}; that classification task differs from reference-relative CO estimation. Spinelle et al. examined field calibration for ozone/NO$_2$ and later NO/CO/CO$_2$, including reference collocation and temporal sensor behavior \cite{spinelle2015,spinelle2017}. Those field conditions and costs are not reproduced by a chamber replay. Recursive identification with forgetting is a standard framework \cite{ljung1983}, rather than a new correction algorithm here.

\paragraph{Conformal prediction.}
Split conformal methods offer marginal coverage under exchangeability of calibration and test data \cite{angelopoulos2021}. Extensions for distribution shift or sequential data need their own conditions. Ordinary split conformal guarantees do not apply unconditionally to arbitrary chronological drift. A properly specified conformal comparator is future work; descriptive envelope containment is not a substitute for its assumptions.
